% Макро-физика (Максвелл, Ньютон, КМ, …) — тултипы в формулах через макросы.

\newcommand{\NotationSecOperatorsTitle}{\textit{Операторы}}
\newcommand{\NotationSecMacroPhysicsTitle}{\textit{Макроскопическая физика (поля, константы, SI)}}
\newcommand{\NotationDescNabla}{векторный дифференциальный оператор
$(\partial/\partial x,\,\partial/\partial y,\,\partial/\partial z)$;
на скаляр $\varphi$ даёт \emph{градиент} $\nabla\varphi$.}
\newcommand{\NotationDescDiv}{дивергенция векторного поля $\mathbf{A}$; $\nabla\cdot\mathbf{A}$
(уравнение непрерывности, Гаусса).}
\newcommand{\NotationDescCurl}{ротор векторного поля $\mathbf{A}$; $\nabla\times\mathbf{A}$
(уравнения Фарадея и Ампера--Максвелла).}
\newcommand{\NotationDescLaplacian}{оператор Лапласа $\nabla^2=\nabla\cdot\nabla$
(уравнение Пуассона, член $-\frac{\hbar^2}{2m}\nabla^2$ в КМ).}
\newcommand{\NotationDescPartialT}{производная по времени $\partial/\partial t$ (локально на $T$);
размерность задаётся операндом.}
\newcommand{\NotationDescEfld}{вектор электрического поля.}
\newcommand{\NotationDescBfld}{вектор магнитной индукции.}
\newcommand{\NotationDescJfld}{плотность электрического тока.}
\newcommand{\NotationDescRhoQ}{объёмная плотность заряда (уравнение Гаусса).}
\newcommand{\NotationDescRhoM}{массовая плотность (гидродинамика, не заряд).}
\newcommand{\NotationDescEpsZero}{электрическая постоянная вакуума \(\varepsilon_0\).}
\newcommand{\NotationDescMuZero}{магнитная постоянная вакуума \(\mu_0\).}
\newcommand{\NotationDescVfld}{вектор скорости.}
\newcommand{\NotationDescRfld}{радиус-вектор точки.}
\newcommand{\NotationDescFfld}{сила (механика Ньютона).}
\newcommand{\NotationDescPfld}{давление (гидродинамика).}
\newcommand{\NotationDescSentropy}{энтропия.}
\newcommand{\NotationDescTtemp}{температура.}
\newcommand{\NotationDescHham}{гамильтониан / оператор Гамильтона.}
\newcommand{\NotationDescPsi}{волновая функция.}
\newcommand{\NotationDescHbar}{постоянная Планка \(\hbar\).}
\newcommand{\NotationDescEcharge}{элементарный заряд \(e\).}
\newcommand{\NotationDescMe}{масса электрона \(m_e\).}
\newcommand{\NotationDescMmass}{масса тела (Ньютон).}
\newcommand{\NotationDescVarphi}{фаза (Маделунг, \(\psi=\sqrt{\rho}\,e^{i\varphi/\hbar}\)).}
\newcommand{\NotationDescClight}{скорость света в уравнениях поля и механики.}
\newcommand{\NotationDesckB}{постоянная Больцмана \(k_B\).}
\newcommand{\NotationDescGgrav}{гравитационная постоянная \(G\).}
\newcommand{\NotationDescVecP}{вектор импульса \(\mathbf{p}\) (пространственная часть \(p^\mu\)).}

\newnotationsectionT{symb:sec-operators}{macro-00}{\NotationSecOperatorsTitle}
\newnotationentryT{symb:nabla}{macro-01}{\ensuremath{\nabla}}{\NotationDescNabla}{\NotationUnitDash}
\newnotationentryT{symb:div}{macro-02}{\ensuremath{\nabla\cdot}}{\NotationDescDiv}{\NotationUnitDash}
\newnotationentryT{symb:curl}{macro-03}{\ensuremath{\nabla\times}}{\NotationDescCurl}{\NotationUnitDash}
\newnotationentryT{symb:laplacian}{macro-04}{\ensuremath{\nabla^2}}{\NotationDescLaplacian}{\NotationUnitDash}
\newnotationentryT{symb:partialt}{macro-05}{\ensuremath{\partial_t}}{\NotationDescPartialT}{\NotationUnitDash}

\newnotationsectionT{symb:sec-macrophys}{macrophys-00}{\NotationSecMacroPhysicsTitle}
\newnotationentryT{symb:Efld}{macrophys-01}{\ensuremath{\mathbf{E}}}{\NotationDescEfld}{\NotationUnitVoltPerMetre}
\newnotationentryT{symb:Bfld}{macrophys-02}{\ensuremath{\mathbf{B}}}{\NotationDescBfld}{\NotationUnitTesla}
\newnotationentryT{symb:Jfld}{macrophys-03}{\ensuremath{\mathbf{j}}}{\NotationDescJfld}{\NotationUnitAmperePerMetreSq}
\newnotationentryT{symb:rhoQ}{macrophys-04}{\ensuremath{\rho}}{\NotationDescRhoQ}{\NotationUnitCoulombPerMetreCubed}
\newnotationentryT{symb:rhoM}{macrophys-05}{\ensuremath{\rho}}{\NotationDescRhoM}{\NotationUnitKgPerMetreCubed}
\newnotationentryT{symb:eps0}{macrophys-06}{\ensuremath{\varepsilon_0}}{\NotationDescEpsZero}{\NotationUnitFaradPerMetre}
\newnotationentryT{symb:mu0}{macrophys-07}{\ensuremath{\mu_0}}{\NotationDescMuZero}{\NotationUnitHenryPerMetre}
\newnotationentryT{symb:Vfld}{macrophys-08}{\ensuremath{\mathbf{v}}}{\NotationDescVfld}{\NotationUnitMetrePerSecond}
\newnotationentryT{symb:Rfld}{macrophys-09}{\ensuremath{\mathbf{r}}}{\NotationDescRfld}{\NotationUnitMetre}
\newnotationentryT{symb:Ffld}{macrophys-10}{\ensuremath{\mathbf{F}}}{\NotationDescFfld}{\NotationUnitNewton}
\newnotationentryT{symb:Pfld}{macrophys-11}{\ensuremath{p}}{\NotationDescPfld}{\NotationUnitPascal}
\newnotationentryT{symb:Sentropy}{macrophys-12}{\ensuremath{S}}{\NotationDescSentropy}{\NotationUnitJoulePerKelvin}
\newnotationentryT{symb:Ttemp}{macrophys-13}{\ensuremath{T}}{\NotationDescTtemp}{\NotationUnitKelvin}
\newnotationentryT{symb:Hham}{macrophys-14}{\ensuremath{H}}{\NotationDescHham}{\NotationUnitJoule}
\newnotationentryT{symb:Psi}{macrophys-15}{\ensuremath{\psi}}{\NotationDescPsi}{\NotationUnitWavefunction}
\newnotationentryT{symb:hbar}{macrophys-16}{\ensuremath{\hbar}}{\NotationDescHbar}{\NotationUnitJouleSecond}
\newnotationentryT{symb:echar}{macrophys-17}{\ensuremath{e}}{\NotationDescEcharge}{\NotationUnitCoulomb}
\newnotationentryT{symb:me}{macrophys-18}{\ensuremath{m_e}}{\NotationDescMe}{\NotationUnitKilogram}
\newnotationentryT{symb:Mmass}{macrophys-19}{\ensuremath{m}}{\NotationDescMmass}{\NotationUnitKilogram}
\newnotationentryT{symb:varphi}{macrophys-20}{\ensuremath{\varphi}}{\NotationDescVarphi}{\NotationUnitJoule}
\newnotationentryT{symb:clight}{macrophys-21}{\ensuremath{c}}{\NotationDescClight}{\NotationUnitMetrePerSecond}
\newnotationentryT{symb:kB}{macrophys-22}{\ensuremath{k_B}}{\NotationDesckB}{\NotationUnitJoulePerKelvin}
\newnotationentryT{symb:Ggrav}{macrophys-23}{\ensuremath{G}}{\NotationDescGgrav}{\NotationUnitGravConstant}
\newnotationentryT{symb:vecP}{macrophys-24}{\ensuremath{\mathbf{p}}}{\NotationDescVecP}{\NotationUnitKgMetrePerSecond}
